\chapter{轮式底盘运动学与非完整约束：从 Pfaffian 矩阵到可控性}\label{ch:mm-wheeled}

% ════════════════════════════════════════════════════════════════
%  P9 移动操作部首章：轮式底盘的速度层运动学。
%  融入 Tutorial「轮式运动学与 Pfaffian」之全部理论，记号归一到本书 SE(2) 约定，
%  Frobenius/可积分类判据指向 ch:constrained-dynamics 去重，此处给轮式工程实例。
%  教学层遵循 v5.0：动机→反面→历史→理论→陷阱→练习 + 符号表/速查/故障排查。
% ════════════════════════════════════════════════════════════════

\begin{practice}[前置自测]
开始之前，先试答下列问题。若已能流畅作答，可只速读本章表格与速查卡；若卡住，正是本章要为你解决的。
\begin{enumerate}
  \item 一辆后轮驱动、前轮转向的汽车，能否\emph{瞬时}向正侧方平移一厘米？经过一连串机动后，能否到达正侧方一米处？这两个回答为什么不矛盾？
  \item 平面刚体的位形空间 $\SE(2)$ 有几个自由度？给体速度 $(v_x^b,v_y^b,\omega)$，世界系下的 $\dot{x},\dot{y}$ 怎么算？
  \item “约束”在数学上是什么？把质点限制在球面上的约束，与把轮子限制为“只能纯滚动”的约束，本质区别在哪？
  \item 给定一族速度约束 $\mathbf{A}(\mathbf{q})\dot{\mathbf{q}}=\mathbf{0}$，你打算如何判断它能否积分成位置约束 $\boldsymbol{\varphi}(\mathbf{q})=\mathbf{0}$？
  \item 差速底盘“不能横移”，可平行泊车时它分明侧向挪进了车位——这个“凭空多出来的侧向运动”从哪来的？能用一个数学对象精确刻画吗？
\end{enumerate}
\end{practice}

\section*{本章目标}
学完本章，你应当能够：
\begin{enumerate}
  \item \textbf{说清}非完整约束“限制瞬时速度方向、却不限制可达位形”的反直觉本质，并与完整约束严格区分；
  \item \textbf{建立} $\SE(2)$ 上的底盘运动学，在体坐标系与世界坐标系之间精确换算速度；
  \item \textbf{独立推导}单轮圆盘、差速、阿克曼、全向/麦轮四种底盘的 Pfaffian 约束矩阵 $\mathbf{A}(\mathbf{q})\dot{\mathbf{q}}=\mathbf{0}$，并据维度/秩读出瞬时自由度；
  \item \textbf{运用} Frobenius 判据（\cref{thm:cd-frobenius} 的轮式专化）判断约束可积性，并用李括号 $[\mathbf{g}_1,\mathbf{g}_2]$ 给出非完整性的显式证明；
  \item \textbf{解释} Chow--Rashevsky 定理与 LARC 如何保证差速车“绕路可达任意位形”，并用 Ball--Box 定理估计绕路代价 $O(\varepsilon^{1/k})$；
  \item \textbf{把}底盘写成漂移自由的仿射控制系统 $\dot{\mathbf{q}}=\sum_i u_i\mathbf{g}_i(\mathbf{q})$，为运动规划与最优控制铺路；
  \item \textbf{建模}纯滚动失效时的轮胎滑移——纵向滑移比 $\kappa$、侧向滑移角 $\alpha$、Pacejka 魔术公式与摩擦锥，并理解从 $\mathbf{A}\dot{\mathbf{q}}=\mathbf{0}$ 到 $\mathbf{A}\dot{\mathbf{q}}=\mathbf{r}(\kappa,\alpha)$ 的约束残差观点及其对里程计漂移的影响。
\end{enumerate}

\subsection*{本章知识导航}
本章是一条“\emph{把轮子与地面的物理接触，翻译成可计算的速度约束、再追问它在全局意义上限制了什么}”的主线。结构如下：

\begin{center}
\small
\begin{tikzpicture}[
  node distance=4.5mm and 7mm, >=Stealth,
  blk/.style={draw, rounded corners, align=center, font=\small, inner sep=3pt, fill=main!6, draw=main!60!black},
  grp/.style={draw, rounded corners, align=center, font=\small, inner sep=3pt, fill=second!6, draw=second!60!black},
  ln/.style={->, thick, gray!60!black}]
\node[blk] (mot) {动机\\不能横移?};
\node[blk, right=of mot] (se2) {$\SE(2)$ 运动学\\体$\leftrightarrow$世界速度};
\node[blk, right=of se2] (pf) {Pfaffian 矩阵\\四种底盘};
\node[grp, below=8mm of mot] (frob) {可积性\\Frobenius};
\node[grp, right=of frob] (chow) {可控性\\Chow / LARC};
\node[grp, right=of chow] (aff) {仿射系统\\$\dot{\mathbf{q}}=\sum u_i\mathbf{g}_i$};
\node[blk, below=8mm of frob, fill=third!8, draw=third!60!black] (slip) {滑移建模\\$\kappa,\alpha$, Pacejka};
\node[blk, right=of slip, fill=third!8, draw=third!60!black] (res) {约束残差\\$\mathbf{A}\dot{\mathbf{q}}=\mathbf{r}$};
\draw[ln] (mot)--(se2); \draw[ln] (se2)--(pf);
\draw[ln] (pf.south) -- ++(0,-4.5mm) -| (frob.north);
\draw[ln] (frob)--(chow); \draw[ln] (chow)--(aff);
\draw[ln] (frob.south)-- (slip.north);
\draw[ln] (slip)--(res);
\end{tikzpicture}
\end{center}

\noindent\textbf{推荐路径}：\cref{sec:mw-motivation,sec:mw-se2} 是任何读者的主干（把“横向不滑”这一句话精确化为速度约束）；\cref{sec:mw-pfaffian} 是核心推导，四种底盘的 $\mathbf{A}(\mathbf{q})$ 一次到位；\cref{sec:mw-integrability,sec:mw-controllability} 是\emph{理论高地}，回答“这约束是否可积、系统是否可控”，并复用约束力学章（\cref{ch:constrained-dynamics}）的判据；\cref{sec:mw-affine} 把底盘写成控制系统的标准形，是后续规划/最优控制章的接口；\cref{sec:mw-slip} 面向真实地面，松弛纯滚动假设。带工程色彩的滑移在线估计与里程计影响可在读完状态估计章后回看。

\subsection*{前置知识桥接}
本章把两条线收束到轮式底盘上。其一是\textbf{约束的分类与判据}：\cref{ch:constrained-dynamics} 已把一切一阶约束统一为 Pfaffian 形式 $\mathbf{A}(\mathbf{q})\dot{\mathbf{q}}+\mathbf{b}=\mathbf{0}$（\cref{eq:cd-pfaffian}），并给出可积性的 Frobenius 判据（\cref{thm:cd-frobenius}）与完整/非完整的形式定义（\cref{def:gm-constraint}）。本章不重复这些一般理论，而是把它们落到四种底盘的\emph{具体矩阵}上，并补全约束力学章一笔带过的可控性侧（Chow 定理）。其二是\textbf{平面位姿的表示}：刚体在平面上的位姿构成李群 $\SE(2)$，是 $\SE(3)$（\cref{ch:lie}）在二维的特例——只有一个转角 $\theta$，于是 $n_q=n_v=3$，本章直接采用这一简化。

\subsection*{如果跳过本章会怎样}
\begin{itemize}
  \item 在运动规划里，你会习惯性地用 A$^\star$/Dijkstra 在栅格上搜最短路，再让差速车跟踪——结果规划出含直角拐弯的路径，底层根本执行不了；不懂\emph{非完整约束}就不知道该换成 Hybrid A$^\star$ 或 state-lattice 这类\emph{尊重运动学}的规划器。
  \item 在轮式里程计里，你会默认轮速就是地面速度；一旦上了湿滑或松软地面，纯滚动假设失效而你毫无察觉，定位会以每秒数厘米的速度漂移——不懂\emph{约束残差} $\mathbf{A}\dot{\mathbf{q}}=\mathbf{r}(\kappa,\alpha)$ 就不知道何时、为何要降低对轮速观测的信任。
  \item 进入移动操作（本部后续章）后，底盘的 $\SE(2)$ 速度要与机械臂雅可比拼成联合雅可比 $\mathbf{V}_{ee}=[\mathbf{J}_b\ \mathbf{J}_a]\boldsymbol{\nu}$；底盘那一块 $\mathbf{J}_b$ 的结构正由本章的 $\mathbf{A}(\mathbf{q})$ 零空间决定。
\end{itemize}

\begin{note}
\textbf{本章记号与约定.}\;
底盘位形 $\mathbf{q}=(x,y,\theta)\in\SE(2)$（或在含轮转角/转向角时维数更高，逐处说明），向量粗体小写、矩阵粗体大写。
体坐标系速度记 $\mathbf{v}^b=(v_x^b,v_y^b,\omega)^\top$，$v_x^b$ 沿车体前向、$v_y^b$ 沿左向、$\omega=\dot\theta$ 为偏航角速度；世界系速度记 $\dot{\mathbf{q}}=(\dot x,\dot y,\dot\theta)^\top$。
非完整约束统一写成 Pfaffian 齐次形式 $\mathbf{A}(\mathbf{q})\dot{\mathbf{q}}=\mathbf{0}$（与 \cref{eq:cd-pfaffian} 在 $\mathbf{b}=\mathbf{0}$、定常情形一致）。
约束零空间张成的\emph{允许速度分布}记 $\mathcal{D}(\mathbf{q})=\ker\mathbf{A}(\mathbf{q})$，其上一组基向量场记 $\mathbf{g}_i(\mathbf{q})$，控制输入记 $u_i$。
李括号 $[\mathbf{g}_i,\mathbf{g}_j]$ 沿用 \cref{ch:lie} 的定义（向量场的对易子）。
\end{note}

% ════════════════════════════════════════════════════════════════
\section{动机：为什么轮式机器人不能随意横移}
\label{sec:mw-motivation}

\paragraph{从停车场说起。}
设想你在驾驶一辆普通轿车（后轮驱动、前轮转向），要从当前车道横移到相邻车道。直觉立刻告诉你：\emph{不能直接横挪}——车子只会前进、后退、转弯。可经过一串“前进—打方向—后退—回正”的操作，你确实能挪进相邻车位。这个再日常不过的经验，藏着一个深刻的数学性质：\textbf{速度受限，位置却不受限}。

\begin{table}[H]\centering\small
\caption{平行泊车经验揭示的三个层次}
\begin{tabular}{ll}
\toprule
特征 & 描述 \\
\midrule
瞬时限制 & 任意时刻速度方向受约束——不能垂直于车轮平面运动 \\
全局自由 & 给足时间与空间，车辆可到达平面上任意位置与朝向 \\
表象矛盾 & 速度受限，可位置不受限——这怎么可能同时成立？ \\
\bottomrule
\end{tabular}
\end{table}

\noindent 这种“瞬时速度受限、可达位形不变”的约束，叫\emph{非完整约束}（nonholonomic constraint）；与之相对的，把系统困在低维曲面上的，叫\emph{完整约束}（holonomic constraint）。二者的形式定义已由 \cref{def:gm-constraint} 给出，这里先把直觉立住。

\paragraph{反面：把两种约束混为一谈会出什么事。}
假设你为差速机器人写路径规划器。若\emph{错把非完整约束当完整约束}：规划器会以为机器人能沿任意方向移动，于是生成含横移、含直角拐弯的轨迹；底层控制器无法执行，机器人原地打转、轮胎刺耳摩擦。反过来，若\emph{错把非完整约束当成不可逾越的禁锢}（以为差速车只能走直线）：大量本可达的位形被误判为不可达，机器人被困在狭窄走廊里直行，丧失了“多次机动到达任意位形”的能力。正确的理解只有一句：非完整约束限制的是\textbf{瞬时运动方向}（微分约束），不是\textbf{最终可达集}（积分约束）。

\paragraph{历史：从经典力学到非线性控制。}
非完整约束的研究始于 $19$ 世纪末的分析力学。Hertz 最早系统讨论它在力学中的角色，Appell 给出处理非完整系统的方程，Chaplygin 求出滚动圆盘等系统的精确解；$1939$ 年 Chow（与 $1938$ 年的 Rashevsky 独立地）证明了非完整系统的可控性判据；$20$ 世纪 $70$ 年代起 Brockett 把非完整约束接入控制论，开创非线性控制；$1994$ 年 Murray、Li 与 Sastry 的《A Mathematical Introduction to Robotic Manipulation》\cite{murray1994mathematical} 把非完整运动规划系统化；其后 LaValle\cite{lavalle2006planning} 又把 RRT 等采样方法推广到非完整系统。本章正是这条脉络在轮式底盘上的落点。

\begin{insight}[完整 vs 非完整的一句话区分]\label{ins:mw-holo}
完整约束减少系统的\textbf{位形空间维度}（质点从三维空间被压到二维球面）；非完整约束减少系统的\textbf{瞬时速度自由度}，却\textbf{不减少}位形空间维度。前者让位形空间“塌缩”，后者只是给每一步的运动方向上了把锁，门后的房间一间不少。
\end{insight}

\begin{pitfall}[概念误区：“非完整 = 不可控”]
新手常想：“轮子不能横移，所以差速车的横向位置不可控。”错。非完整约束锁的是\emph{瞬时速度方向}，不是\emph{可达位形集}。差速车靠“前进—转弯—前进”的序列能到达平面上任意 $(x,y,\theta)$；其可控性由 Chow 定理保证（\cref{sec:mw-controllability}）。“不可控”指存在永远到不了的状态；“非完整”只是说到某些状态要\emph{绕路}。务必把这两件事分开。
\end{pitfall}

\begin{practice}
\begin{enumerate}
  \item 各举三例日常的完整约束与非完整约束，并说明每个限制的是位置还是速度。
  \item 一个冰球在冰面上自由滑行，它受的约束（若有）是完整还是非完整？若再加上“只能沿当前朝向滑”的约束，情况如何变化？
  \item 平面刚体位形空间 $\SE(2)$ 维度为 $3$。若施加一个非完整约束减少一个瞬时速度自由度，剩几个瞬时速度自由度？位形空间维度变了吗？
\end{enumerate}
\end{practice}

% ════════════════════════════════════════════════════════════════
\section{\texorpdfstring{$\SE(2)$}{SE(2)} 上的运动学与速度约束}
\label{sec:mw-se2}

\paragraph{位形与速度的两套坐标。}
轮式底盘在平面上运动，位形完全由\emph{位置} $(x,y)$ 与\emph{朝向} $\theta$ 确定，构成李群
\begin{equation}
  \mathbf{q}=\begin{pmatrix}x\\y\\\theta\end{pmatrix}\in\SE(2)\cong\mathbb{R}^2\times S^1,
  \qquad
  \mathbf{T}=\begin{pmatrix}\cos\theta&-\sin\theta&x\\\sin\theta&\cos\theta&y\\0&0&1\end{pmatrix}\in\SE(2).
  \label{eq:mw-se2}
\end{equation}
与 $\SE(3)$（\cref{ch:lie}）相比，平面上朝向只剩一个角 $\theta$，无需四元数，故位形与速度同维：$n_q=n_v=3$。这是 $S^1$ 相对 $\SO(3)$ 的简化，也是后续推导都在 $3$ 维里打转的原因。

\paragraph{体速度与世界速度的换算。}
刚体速度可在两个坐标系里写。\emph{世界系速度}（空间速度）就是位形的时间导数 $\dot{\mathbf{q}}=(\dot x,\dot y,\dot\theta)^\top$；\emph{体速度} $\mathbf{v}^b=(v_x^b,v_y^b,\omega)^\top$ 则把线速度投影到随车转动的前向/左向轴上。二者经平面旋转矩阵相联：
\begin{equation}
  \begin{pmatrix}\dot x\\\dot y\end{pmatrix}
  =\mathbf{R}(\theta)\begin{pmatrix}v_x^b\\v_y^b\end{pmatrix}
  =\begin{pmatrix}\cos\theta&-\sin\theta\\\sin\theta&\cos\theta\end{pmatrix}\begin{pmatrix}v_x^b\\v_y^b\end{pmatrix},
  \qquad \dot\theta=\omega.
  \label{eq:mw-body2world}
\end{equation}
逐分量展开即
\begin{equation}
  \dot x=v_x^b\cos\theta-v_y^b\sin\theta,\quad
  \dot y=v_x^b\sin\theta+v_y^b\cos\theta,\quad
  \dot\theta=\omega.
  \label{eq:mw-body2world-comp}
\end{equation}
这组式子是本章一切 Pfaffian 推导的“坐标基建”：物理约束写在体系里最自然（“侧向速度为零”就是 $v_y^b=0$），而约束矩阵 $\mathbf{A}(\mathbf{q})$ 必须落到世界系的 $\dot{\mathbf{q}}$ 上，二者靠 \cref{eq:mw-body2world} 来回搬运。

\begin{figure}[ht]
\centering
\begin{tikzpicture}[>=Stealth, scale=1.0]
  % world frame
  \draw[->, gray!70!black] (-0.4,-0.4) -- (1.4,-0.4) node[right]{$x$};
  \draw[->, gray!70!black] (-0.4,-0.4) -- (-0.4,1.2) node[above]{$y$};
  \node[gray!70!black] at (-0.7,-0.7) {世界系};
  % chassis at (2.6,0.9), heading theta
  \begin{scope}[shift={(3.0,0.7)}, rotate=28]
    \draw[thick, main!70!black, fill=main!8, rounded corners=2pt] (-0.7,-0.45) rectangle (0.7,0.45);
    % wheels
    \fill[gray!50!black] (-0.55,-0.5) rectangle (-0.25,-0.4);
    \fill[gray!50!black] (-0.55,0.4) rectangle (-0.25,0.5);
    \fill[gray!50!black] (0.25,-0.5) rectangle (0.55,-0.4);
    \fill[gray!50!black] (0.25,0.4) rectangle (0.55,0.5);
    % body axes
    \draw[->, second!70!black, very thick] (0,0) -- (1.1,0) node[right]{$\mathbf{e}_\parallel$ (前向 $v_x^b$)};
    \draw[->, third!70!black, very thick] (0,0) -- (0,0.95) node[above]{$\mathbf{e}_\perp$ ($v_y^b$)};
    \fill (0,0) circle (1.2pt);
  \end{scope}
  % position + heading annotation
  \draw[->, dashed, gray!60!black] (-0.4,-0.4) -- (3.0,0.7);
  \node[gray!60!black] at (1.3,0.45) {$(x,y)$};
  \draw[->, gray!70!black] (3.7,0.7) arc (0:28:0.7);
  \node[gray!70!black] at (4.1,0.95) {$\theta$};
\end{tikzpicture}
\caption{$\SE(2)$ 底盘的体坐标系与世界坐标系。前向轴 $\mathbf{e}_\parallel$、左向轴 $\mathbf{e}_\perp$ 随底盘转过 $\theta$；体速度 $(v_x^b,v_y^b,\omega)$ 经 \cref{eq:mw-body2world} 变到世界系 $(\dot x,\dot y,\dot\theta)$。}
\label{fig:mw-se2-frames}
\end{figure}

\begin{insight}[反事实：若不分体系与世界系会怎样]\label{ins:mw-frame}
取一辆朝向 $\theta=90^\circ$ 的车，以体系前向速度 $v_x^b=1$\,m/s 行驶。世界系里它其实在沿 $y$ 方向走（$\dot x=0,\dot y=1$）。若把体速度直接当世界速度用，就会算出车在沿 $x$ 走——方向完全错。这类错误在仿真里表现为机器人“螃蟹横行”。\cref{eq:mw-body2world} 的旋转矩阵正是杜绝它的唯一闸门。
\end{insight}

\paragraph{Pfaffian 形式的精确定义。}
现在可以把约束精确化。\emph{Pfaffian 约束}是对广义速度的线性约束，本章取其定常齐次形式
\begin{equation}
  \mathbf{A}(\mathbf{q})\dot{\mathbf{q}}=\mathbf{0},\qquad \mathbf{A}(\mathbf{q})\in\mathbb{R}^{k\times n},
  \label{eq:mw-pfaffian}
\end{equation}
其中 $k$ 是独立约束个数、$n$ 是位形维度。这正是 \cref{eq:cd-pfaffian} 在 $\mathbf{b}=\mathbf{0}$ 时的特例。$\mathbf{A}(\mathbf{q})$ 的每一行指定一个\emph{被禁止的速度方向}：系统速度不得有沿该方向的分量。所有合规速度构成\emph{允许速度分布}
\begin{equation}
  \mathcal{D}(\mathbf{q})=\{\mathbf{v}\in\mathbb{R}^n\mid\mathbf{A}(\mathbf{q})\mathbf{v}=\mathbf{0}\}=\ker\mathbf{A}(\mathbf{q}),
  \label{eq:mw-distribution}
\end{equation}
其维度 $\dim\mathcal{D}=n-\operatorname{rank}\mathbf{A}(\mathbf{q})$，即系统在该时刻的\textbf{瞬时运动自由度}。

\begin{table}[H]\centering\small
\caption{Pfaffian 约束的四个关键量}
\begin{tabular}{lll}
\toprule
概念 & 数学表示 & 物理含义 \\
\midrule
位形空间维度 & $n=\dim(\mathbf{q})$ & 描述位置需几个坐标 \\
约束个数 & $k=\operatorname{rank}(\mathbf{A})$ & 几个独立速度方向被禁 \\
瞬时自由度 & $n-k$ & 此刻能沿几个方向运动 \\
允许速度分布 & $\mathcal{D}(\mathbf{q})=\ker\mathbf{A}(\mathbf{q})$ & 此刻所有合规速度的集合 \\
\bottomrule
\end{tabular}
\end{table}

\begin{pitfall}[编程陷阱：在世界系写约束却忘了它依赖 $\theta$]
想表达“不能横移”，新手常直接写常数矩阵 $\mathbf{A}=[0,1,0]$（禁 $y$ 方向速度）。这只在 $\theta=0$ 时对，其他朝向全错。根因是 $\mathbf{A}(\mathbf{q})$ \emph{必须}是位形（尤其 $\theta$）的函数：“禁横向速度”在体系是常数约束 $v_y^b=0$，但经 \cref{eq:mw-body2world-comp} 搬到世界系就成了 $-\dot x\sin\theta+\dot y\cos\theta=0$，矩阵 $\mathbf{A}$ 随 $\theta$ 转。正确做法永远是：先在体系写约束，再用旋转矩阵变到世界系。
\end{pitfall}

\begin{insight}[轮式约束与足端接触是同一数学对象]\label{ins:mw-unify}
足式机器人的“足端静止接触” $\mathbf{J}_c(\mathbf{q})\mathbf{v}=\mathbf{0}$（见统一动力学 \cref{eq:umm-core} 中的接触项）与轮式的 $\mathbf{A}(\mathbf{q})\dot{\mathbf{q}}=\mathbf{0}$ \emph{形式完全相同}——令 $\mathbf{A}=\mathbf{J}_c$ 即可。区别只在物理：足端约束要求接触点速度为零（完整约束，可积成“接触点位置固定”），轮式约束要求接触点不横滑（非完整约束，不可积成位置关系）。这一统一视角是把底盘嵌入浮基统一动力学、乃至轮足混合系统的数学基础。
\end{insight}

\begin{practice}
\begin{enumerate}
  \item 写出 $\SE(2)$ 上从体速度到世界速度的完整 $3\times3$ 变换矩阵，并验证其行列式为 $1$（变换保体积）。
  \item 把约束 $v_y^b=0$ 转成 Pfaffian 形式 $\mathbf{A}(\mathbf{q})\dot{\mathbf{q}}=0$，写出 $\mathbf{A}(\mathbf{q})$，并核对 $\theta=0,\pi/4,\pi/2$ 时的具体形态。
  \item 一个机器人受约束 $v_x^b=v_y^b=0$（只能原地转）。这是完整还是非完整约束？允许速度分布维度是多少？
\end{enumerate}
\end{practice}

% ════════════════════════════════════════════════════════════════
\section{四种底盘的 Pfaffian 约束矩阵}
\label{sec:mw-pfaffian}

本节是全章核心。我们从物理约束出发，逐一推导四种主流底盘的 Pfaffian 矩阵。每种推导都走同一个\textbf{三步框架}：(1) 写出体系物理约束（纯滚动 + 不横滑）；(2) 经 \cref{eq:mw-body2world-comp} 变到世界系；(3) 组装 $\mathbf{A}(\mathbf{q})$，并据秩读出瞬时自由度。

\subsection{单轮圆盘：最小非完整系统}
\label{sec:mw-disk}

竖直圆盘在平面上滚动，是非完整约束的\emph{最小模型}：四个位形坐标、两个约束。它的配置由轮心位置、轮平面朝向、绕轴转角描述：
\begin{equation}
  \mathbf{q}=(x,\ y,\ \varphi,\ \theta)^\top\in\mathbb{R}^2\times S^1\times S^1,\qquad n=4,
  \label{eq:mw-disk-q}
\end{equation}
其中 $\varphi$ 是轮平面朝向（heading），$\theta$ 是绕自身轴的滚动角。在轮的体系里定义前向单位向量 $\mathbf{e}_\parallel=(\cos\varphi,\sin\varphi)^\top$ 与横向单位向量 $\mathbf{e}_\perp=(-\sin\varphi,\cos\varphi)^\top$。

\begin{derivation}[单轮圆盘的两条纯滚动约束]
记轮心平面位置 $\mathbf{p}=(x,y)^\top$。\textbf{约束一（纵向纯滚动）}：轮心沿前向的速度等于轮缘线速度 $r\dot\theta$——轮子向前滚一圈，轮心走过的距离等于轮缘展开的弧长。若不满足，轮子要么在地上空转、要么被拖着滑：
\begin{equation}
  \mathbf{e}_\parallel^\top\dot{\mathbf{p}}-r\dot\theta=0
  \;\Longrightarrow\;
  \dot x\cos\varphi+\dot y\sin\varphi-r\dot\theta=0.
  \label{eq:mw-disk-roll}
\end{equation}
\textbf{约束二（横向不滑）}：轮心沿轮轴方向的速度为零——轮子不侧滑（否则就像车在冰面上打横）：
\begin{equation}
  \mathbf{e}_\perp^\top\dot{\mathbf{p}}=0
  \;\Longrightarrow\;
  -\dot x\sin\varphi+\dot y\cos\varphi=0.
  \label{eq:mw-disk-noslip}
\end{equation}
把两式写成矩阵形式 $\mathbf{A}(\mathbf{q})\dot{\mathbf{q}}=\mathbf{0}$：
\begin{equation}
  \underbrace{\begin{pmatrix}\cos\varphi&\sin\varphi&0&-r\\-\sin\varphi&\cos\varphi&0&0\end{pmatrix}}_{\mathbf{A}(\mathbf{q})}
  \begin{pmatrix}\dot x\\\dot y\\\dot\varphi\\\dot\theta\end{pmatrix}=\begin{pmatrix}0\\0\end{pmatrix}.
  \label{eq:mw-disk-A}
\end{equation}
\end{derivation}

\noindent\textbf{维度核验}：$\mathbf{A}\in\mathbb{R}^{2\times4}$，两行线性无关（前两列构成的 $2\times2$ 块行列式 $\cos^2\varphi+\sin^2\varphi=1\neq0$），故 $\operatorname{rank}\mathbf{A}=2$，瞬时自由度 $=4-2=2$。物理上：圆盘任意时刻只能做两种运动——前/后滚（同时改变 $\theta$ 与 $(x,y)$）和改变朝向（改变 $\varphi$）；既不能横移，也不能在不转轮的情况下平移。

求 \cref{eq:mw-disk-A} 的零空间，得允许速度的参数化
\begin{equation}
  \dot{\mathbf{q}}=u_1\underbrace{\begin{pmatrix}r\cos\varphi\\r\sin\varphi\\0\\1\end{pmatrix}}_{\mathbf{g}_1(\mathbf{q})}
  +u_2\underbrace{\begin{pmatrix}0\\0\\1\\0\end{pmatrix}}_{\mathbf{g}_2(\mathbf{q})},
  \label{eq:mw-disk-vf}
\end{equation}
其中 $u_1=\dot\theta$（轮转速）、$u_2=\dot\varphi$（转向速率）是两个输入，$\mathbf{g}_1,\mathbf{g}_2$ 是张成 $\mathcal{D}(\mathbf{q})$ 的两个\emph{控制向量场}。

\begin{insight}[Pfaffian 零空间即控制向量场]\label{ins:mw-vf}
约束 $\mathbf{A}(\mathbf{q})\dot{\mathbf{q}}=\mathbf{0}$ 的零空间给出一组控制向量场 $\{\mathbf{g}_1,\dots,\mathbf{g}_m\}$，系统运动方程随之写成 $\dot{\mathbf{q}}=\sum_i u_i\mathbf{g}_i(\mathbf{q})$——这是\emph{漂移自由（drift-free）的仿射控制系统}的标准形（\cref{sec:mw-affine}）。非完整运动规划的全部工具（Chow 定理、RRT、steering function）都建立在这个形式之上。
\end{insight}

\subsection{差速底盘}
\label{sec:mw-diff}

差速底盘（differential drive）是最简单的实用底盘：两个独立驱动轮加一或多个被动万向轮，TurtleBot、iRobot Create 与大量教学机器人皆是。其位形只需三个变量（轮转角由运动学约束隐含）：$\mathbf{q}=(x,y,\theta)^\top\in\SE(2)$，$(x,y)$ 取两轮中点、$\theta$ 为底盘朝向。两轮转速 $\omega_L,\omega_R$ 映射到底盘速度：
\begin{equation}
  v_x^b=\frac{r(\omega_R+\omega_L)}{2},\qquad \omega=\dot\theta=\frac{r(\omega_R-\omega_L)}{L},
  \label{eq:mw-diff-fk}
\end{equation}
$r$ 为轮半径、$L$ 为轮距。差速底盘的\emph{唯一}约束是横向不滑 $v_y^b=0$，经 \cref{eq:mw-body2world-comp} 变到世界系：
\begin{equation}
  -\dot x\sin\theta+\dot y\cos\theta=0
  \;\Longrightarrow\;
  \mathbf{A}(\mathbf{q})=\begin{pmatrix}-\sin\theta&\cos\theta&0\end{pmatrix}.
  \label{eq:mw-diff-A}
\end{equation}
维度分析：$n=3$，$k=\operatorname{rank}\mathbf{A}=1$，瞬时自由度 $=3-1=2$（前进/后退 $+$ 原地旋转）。注意它与单轮圆盘的关键差别：差速底盘有两个独立电机输入 $(\omega_L,\omega_R)$，恰好等于两个瞬时自由度——系统在允许速度子空间内是\emph{完全驱动}的。运动方程
\begin{equation}
  \dot{\mathbf{q}}=\underbrace{\begin{pmatrix}\cos\theta\\\sin\theta\\0\end{pmatrix}}_{\mathbf{g}_1}v
  +\underbrace{\begin{pmatrix}0\\0\\1\end{pmatrix}}_{\mathbf{g}_2}\omega,
  \label{eq:mw-diff-vf}
\end{equation}
$v=v_x^b$ 为前向线速度、$\omega=\dot\theta$ 为角速度。这两个向量场将在 \cref{sec:mw-integrability,sec:mw-controllability} 反复出场，是全章的“主角”。

\subsection{阿克曼转向}
\label{sec:mw-ackermann}

阿克曼转向（Ackermann steering）是汽车与类车机器人的标准转向方式：转弯靠前轮偏转，而非两轮差速。其位形需四个变量：
\begin{equation}
  \mathbf{q}=(x,\ y,\ \theta,\ \delta)^\top,
  \label{eq:mw-ack-q}
\end{equation}
$(x,y)$ 取后轴中点、$\theta$ 为车体朝向、$\delta$ 为前轮转向角。约束仍是后轮不横滑：
\begin{equation}
  v_y^b=0\;\Longrightarrow\;-\dot x\sin\theta+\dot y\cos\theta=0,\qquad
  \mathbf{A}(\mathbf{q})=\begin{pmatrix}-\sin\theta&\cos\theta&0&0\end{pmatrix}.
  \label{eq:mw-ack-A}
\end{equation}
阿克曼运动学额外给出转向角 $\delta$ 与转弯半径 $R$ 的几何关系 $\tan\delta=L_{\text{wb}}/R$（$L_{\text{wb}}$ 为轴距），从而把偏航角速度与前向速度焊在一起：
\begin{equation}
  \dot\theta=\frac{v}{R}=\frac{v\tan\delta}{L_{\text{wb}}}.
  \label{eq:mw-ack-yaw}
\end{equation}
维度分析与差速相同（$n=4$ 时 $k=1$，对 $(x,y,\theta)$ 子空间瞬时自由度仍为这套约束所限），但实际输入只有 $2$ 个（前向速度 $v$、转向角速率 $\dot\delta$），故系统是\emph{欠驱动}的——$\delta$ 不受 Pfaffian 约束限制，却受机械转向角上限 $|\delta|\le\delta_{\max}$ 约束。运动方程
\begin{equation}
  \dot{\mathbf{q}}=\begin{pmatrix}\cos\theta\\\sin\theta\\\tan\delta/L_{\text{wb}}\\0\end{pmatrix}v
  +\begin{pmatrix}0\\0\\0\\1\end{pmatrix}\dot\delta.
  \label{eq:mw-ack-vf}
\end{equation}

\begin{insight}[反事实：若阿克曼没有最大转向角]\label{ins:mw-ackmax}
设想 $\delta_{\max}=90^\circ$。当 $\delta\to90^\circ$ 时 $\tan\delta\to\infty$，即便前向速度 $v$ 很小，偏航角速度 $\dot\theta$ 也趋于无穷——车辆绕后轴做极小半径回转，物理上退化为差速车的原地旋转。真实汽车 $\delta_{\max}$ 通常只有 $30^\circ$–$40^\circ$，这进一步抬高了最小转弯半径，正是停车入库要反复“揉库”的根源。
\end{insight}

\subsection{全向轮与麦克纳姆轮}
\label{sec:mw-omni}

全向底盘（如 Kuka youBot 底座、RoboCup 小型组机器人）用特殊轮消去横向不滑约束，使机器人能沿任意方向瞬时运动。\emph{全向轮}（Swedish/omni wheel）在轮缘装一圈自由滚动的小辊子，使轮子可沿轴向自由滑动，于是横向不滑约束被\emph{消除}。三轮全向底盘（$120^\circ$ 对称布置）的运动方程为
\begin{equation}
  \dot{\mathbf{q}}=\mathbf{J}^{-1}(\theta)\,(r\omega_1,\ r\omega_2,\ r\omega_3)^\top,
  \label{eq:mw-omni-fk}
\end{equation}
而其 Pfaffian 约束矩阵 $\mathbf{A}(\mathbf{q})=\varnothing$——\textbf{无约束}，全向底盘是\emph{完整系统}，位形空间维度 $=$ 瞬时自由度 $=3$，可瞬时沿任意方向（前后、左右、旋转的任意组合）运动。

\emph{麦克纳姆轮}（Mecanum wheel）是全向运动的另一实现：辊子以 $45^\circ$ 角安装（全向轮是 $90^\circ$）。四轮麦轮底盘（矩形布置）的逆运动学为
\begin{equation}
  \begin{pmatrix}\omega_1\\\omega_2\\\omega_3\\\omega_4\end{pmatrix}
  =\frac{1}{r}\begin{pmatrix}1&-1&-(L_x+L_y)\\1&1&(L_x+L_y)\\1&1&-(L_x+L_y)\\1&-1&(L_x+L_y)\end{pmatrix}
  \begin{pmatrix}v_x^b\\v_y^b\\\omega\end{pmatrix},
  \label{eq:mw-mecanum}
\end{equation}
$L_x,L_y$ 为轮到底盘中心的前后/左右距离。与全向轮一样，麦轮底盘\emph{没有}非完整约束，是完整系统；但工程上辊子与地面的接触力学使其\textbf{横向力传递效率低}，故侧向运动的最大加速度远小于前向——这是\emph{动力学}限制，而非运动学约束，切勿混入 $\mathbf{A}(\mathbf{q})$。

\begin{table}[H]\centering\small
\caption{四种底盘的 Pfaffian 矩阵与瞬时自由度对照}
\label{tab:mw-chassis}
\begin{tabular}{lccc c}
\toprule
底盘类型 & $\mathbf{A}(\mathbf{q})$ & $n$ & $\operatorname{rank}\mathbf{A}$ & 瞬时自由度 / 非完整？ \\
\midrule
单轮圆盘 & $\left[\begin{smallmatrix}\cos\varphi&\sin\varphi&0&-r\\-\sin\varphi&\cos\varphi&0&0\end{smallmatrix}\right]$ & $4$ & $2$ & $2$ / 是 \\[2pt]
差速 & $\left[\begin{smallmatrix}-\sin\theta&\cos\theta&0\end{smallmatrix}\right]$ & $3$ & $1$ & $2$ / 是 \\[2pt]
阿克曼 & $\left[\begin{smallmatrix}-\sin\theta&\cos\theta&0&0\end{smallmatrix}\right]$ & $4$ & $1$ & $3$（输入 $2$）/ 是 \\[2pt]
全向 / 麦轮 & $\varnothing$ & $3$ & $0$ & $3$ / 否 \\
\bottomrule
\end{tabular}
\end{table}

\begin{pitfall}[思维陷阱：多轮约束当成各轮约束的简单叠加]
四轮系统看似每个轮各两个约束、共 $8$ 个。但 $\SE(2)$ 只有 $3$ 个自由度，这 $8$ 个约束至多 $1$–$3$ 个独立，其余线性相关——是哪些独立，由轮距、轴距等\emph{轮间几何}决定。正确做法：先写出所有轮的约束，再分析堆叠后 $\mathbf{A}(\mathbf{q})$ 的秩，只保留独立行。直接“每个轮各自满足约束就行”会得到一个秩亏、自相矛盾的约束系统。
\end{pitfall}

\begin{pitfall}[概念误区：“全向底盘无约束，所以规划最简单”]
全向底盘\emph{无运动学约束}，但仍有\emph{动力学约束}（最大加速度、轮打滑上限）与\emph{环境约束}（避障）。况且其横向力传递效率低，实际最大横向加速度往往只有前向的 $50\%$–$70\%$。运动规划的难度不只来自运动学约束，还来自动力学与环境约束——“能横移”不等于“随便画直线就能跟踪”。
\end{pitfall}

\begin{practice}
\begin{enumerate}
  \item 为差速底盘推导正运动学：给 $\omega_L,\omega_R$，算 $(\dot x,\dot y,\dot\theta)$；验证 $\omega_L=\omega_R$ 时直行、$\omega_L=-\omega_R$ 时原地旋转。
  \item 推导自行车模型（阿克曼的简化，$\mathbf{q}=(x,y,\theta,\delta)$，$(x,y)$ 取后轮接地点）的 Pfaffian 约束矩阵，与汽车阿克曼矩阵 \cref{eq:mw-ack-A} 对比。
  \item 用符号计算（如 SymPy）求四轮麦轮底盘的正运动学矩阵，验证四轮同速同向时前进、对角轮反向时原地旋转。
\end{enumerate}
\end{practice}

% ════════════════════════════════════════════════════════════════
\section{可积性：Frobenius 判据的轮式专化}
\label{sec:mw-integrability}

\paragraph{悬而未决的问题。}
上一节推了四种底盘的 Pfaffian 约束。一个关键问题仍未回答：\textbf{如何判断一个 Pfaffian 约束是非完整的？}也就是说，$\mathbf{A}(\mathbf{q})\dot{\mathbf{q}}=\mathbf{0}$ 能否积分成位置约束 $\boldsymbol{\varphi}(\mathbf{q})=\mathbf{0}$？这不是数学好奇心——它直接决定规划策略：约束若可积，可达空间降维，规划在低维流形上做；约束若不可积，可达空间不变，但每步运动方向受限，规划要“绕路”。判错类型，规划器就会给出执行不了的路径。

一般判据已由约束力学章给出，本节只做轮式专化，不重复证明：

\begin{theorem}[可积判据 = Frobenius（轮式专化，引自 \cref{thm:cd-frobenius}）]\label{thm:mw-frobenius}
设允许速度分布 $\mathcal{D}(\mathbf{q})=\ker\mathbf{A}(\mathbf{q})=\operatorname{span}\{\mathbf{g}_1,\dots,\mathbf{g}_m\}$ 光滑。则 $\mathcal{D}$ 可积（存在 $n-m$ 个独立函数 $\varphi_i$ 使 $\nabla\varphi_i\perp\mathcal{D}$，约束完整）当且仅当 $\mathcal{D}$ \emph{对合}（involutive）：
\begin{equation}
  \forall\,i,j:\quad [\mathbf{g}_i,\mathbf{g}_j](\mathbf{q})\in\mathcal{D}(\mathbf{q}),
  \label{eq:mw-involutive}
\end{equation}
其中 $[\mathbf{g}_i,\mathbf{g}_j]=\dfrac{\partial\mathbf{g}_j}{\partial\mathbf{q}}\mathbf{g}_i-\dfrac{\partial\mathbf{g}_i}{\partial\mathbf{q}}\mathbf{g}_j$ 为李括号。等价地，以 $\omega^i$ 记 $\mathbf{A}$ 的第 $i$ 行（约束 $1$-形式），可积当且仅当 $\mathrm{d}\omega^i\wedge\omega^1\wedge\cdots\wedge\omega^k=0\ (\forall i)$；单约束（$k=1$）时简化为 $\mathrm{d}\omega\wedge\omega=0$。
\end{theorem}

\begin{insight}[李括号的几何含义]\label{ins:mw-bracket}
$[\mathbf{g}_i,\mathbf{g}_j]$ 度量“先沿 $\mathbf{g}_i$ 走一小步、再沿 $\mathbf{g}_j$ 走一小步”与“先 $\mathbf{g}_j$ 再 $\mathbf{g}_i$”两条路径终点之差。若两终点重合（括号为零或落在 $\mathcal{D}$ 内），约束可积；若终点不同（括号逸出 $\mathcal{D}$），说明交替使用两个输入能\emph{生成新的运动方向}——这正是非完整系统“绕路到达”的数学本质。它也类比微积分里的判据：向量场 $\mathbf{F}$ 是某标量梯度当且仅当 $\nabla\times\mathbf{F}=\mathbf{0}$（无旋）；Frobenius 定理把“无旋”推广成“对合”。
\end{insight}

\subsection{差速底盘：李括号显式计算}
\label{sec:mw-diff-bracket}

用 \cref{thm:mw-frobenius} 验证差速底盘的约束确实非完整。两个控制向量场取自 \cref{eq:mw-diff-vf}：
\begin{equation}
  \mathbf{g}_1=(\cos\theta,\ \sin\theta,\ 0)^\top,\qquad \mathbf{g}_2=(0,\ 0,\ 1)^\top.
  \label{eq:mw-diff-g}
\end{equation}

\begin{derivation}[{$[\mathbf{g}_1,\mathbf{g}_2]$ 恰是体系侧向方向}]
两个雅可比为
\[
  \frac{\partial\mathbf{g}_1}{\partial\mathbf{q}}=\begin{pmatrix}0&0&-\sin\theta\\0&0&\cos\theta\\0&0&0\end{pmatrix},\qquad
  \frac{\partial\mathbf{g}_2}{\partial\mathbf{q}}=\mathbf{0}.
\]
于是
\begin{equation}
  [\mathbf{g}_1,\mathbf{g}_2]
  =\frac{\partial\mathbf{g}_2}{\partial\mathbf{q}}\mathbf{g}_1-\frac{\partial\mathbf{g}_1}{\partial\mathbf{q}}\mathbf{g}_2
  =\mathbf{0}-\begin{pmatrix}-\sin\theta\\\cos\theta\\0\end{pmatrix}
  =\begin{pmatrix}\sin\theta\\-\cos\theta\\0\end{pmatrix}.
  \label{eq:mw-diff-bracket}
\end{equation}
检验它是否落在 $\mathcal{D}=\operatorname{span}\{\mathbf{g}_1,\mathbf{g}_2\}$ 内：若存在 $\alpha,\beta$ 使 $(\sin\theta,-\cos\theta,0)^\top=\alpha(\cos\theta,\sin\theta,0)^\top+\beta(0,0,1)^\top$，由第三行得 $\beta=0$，前两行给 $\alpha\cos\theta=\sin\theta$ 且 $\alpha\sin\theta=-\cos\theta$，相除得 $\tan\theta=-1/\tan\theta$，即 $\tan^2\theta=-1$，\emph{无实解}。故 $[\mathbf{g}_1,\mathbf{g}_2]\notin\mathcal{D}$，分布非对合，约束\textbf{不可积}，系统\textbf{非完整}。
\end{derivation}

\noindent 也可走微分形式一路核验：差速约束 $1$-形式 $\omega=-\sin\theta\,\mathrm{d}x+\cos\theta\,\mathrm{d}y$，其外微分 $\mathrm{d}\omega=-\cos\theta\,\mathrm{d}\theta\wedge\mathrm{d}x-\sin\theta\,\mathrm{d}\theta\wedge\mathrm{d}y$，与 $\omega$ 楔乘并用 $\mathrm{d}\theta\wedge\mathrm{d}y\wedge\mathrm{d}x=-\mathrm{d}\theta\wedge\mathrm{d}x\wedge\mathrm{d}y$ 合并，得
\begin{equation}
  \mathrm{d}\omega\wedge\omega=(\cos^2\theta+\sin^2\theta)\,\mathrm{d}\theta\wedge\mathrm{d}x\wedge\mathrm{d}y=\mathrm{d}\theta\wedge\mathrm{d}x\wedge\mathrm{d}y\neq0,
  \label{eq:mw-diff-form}
\end{equation}
两条路殊途同归：约束不可积。这与约束力学章对独轮车的同一结论（\cref{ex:cd-three}）一致。

\begin{insight}[平行泊车的数学本质]\label{ins:mw-parking}
$[\mathbf{g}_1,\mathbf{g}_2]=(\sin\theta,-\cos\theta,0)^\top$ 恰是体系的\emph{侧向}方向。这意味着：交替使用前进（$\mathbf{g}_1$）与转向（$\mathbf{g}_2$），差速车能凑出侧向运动——这正是平行泊车的内核。沿“前进 $\varepsilon$ → 转 $\varepsilon$ → 后退 $\varepsilon$ → 反转 $\varepsilon$”走一圈，机器人不回到原点，而是净侧移 $\Delta\mathbf{q}\approx\varepsilon^2[\mathbf{g}_1,\mathbf{g}_2]+O(\varepsilon^3)$。每轮机动产生 $O(\varepsilon^2)$ 的侧移，故要反复“揉”多次才挪进车位。李括号不是抽象记号，它精确地说出了“靠已有运动方向的交替，能造出哪个新方向”。
\end{insight}

\begin{pitfall}[编程陷阱：数值算李括号时步长选大了]
用有限差分近似 $\partial\mathbf{g}/\partial\mathbf{q}$ 时取步长 $h=0.1$（过大），算出的数值李括号与解析结果差异显著，尤其当 $\theta$ 接近 $0$ 或 $\pi$。根因是 $\mathbf{g}_1=(\cos\theta,\sin\theta,0)^\top$ 对 $\theta$ 的导数是三角函数，有限差分需很小步长才精确。正确做法：用符号微分（SymPy），或自动微分（如 JAX 的前向模式）。
\end{pitfall}

\begin{practice}
\begin{enumerate}
  \item 计算单轮圆盘（\cref{eq:mw-disk-vf} 的两个向量场）的李括号 $[\mathbf{g}_1,\mathbf{g}_2]$，验证它不在 $\mathcal{D}=\operatorname{span}\{\mathbf{g}_1,\mathbf{g}_2\}$ 内。
  \item 系统 $\mathbf{q}=(x,y,\theta)$ 受约束 $\dot x\sin\theta-\dot y\cos\theta+\ell\dot\theta=0$（$\ell$ 常数，拖车约束）。用 Frobenius 判据判断它是否可积。
  \item 证明全向底盘的约束平凡可积（因无约束）；再给全向底盘挂一节拖车（底盘 $(x,y,\theta)$ 无约束、拖车角 $\phi$ 受约束），写出整系统的 Pfaffian 矩阵并判可积性。
\end{enumerate}
\end{practice}

% ════════════════════════════════════════════════════════════════
\section{可控性：Chow--Rashevsky 定理与绕路代价}
\label{sec:mw-controllability}

\paragraph{从“不可积”到“可达”。}
Frobenius 判据告诉我们差速车的约束非完整——瞬时不能横移。但“瞬时不能横移”绝不等于“到不了侧方位置”。Chow--Rashevsky 定理回答最关键的问题：非完整系统到底能到达哪些位形？答案靠李代数秩条件给出。

\begin{theorem}[李代数秩条件（LARC）与 Chow--Rashevsky\cite{murray1994mathematical}]\label{thm:mw-chow}
设漂移自由控制系统 $\dot{\mathbf{q}}=\sum_{i=1}^m u_i\mathbf{g}_i(\mathbf{q})$，$\mathbf{g}_i$ 光滑。若在每一点 $\mathbf{q}$，向量场 $\{\mathbf{g}_i\}$ 及其各阶李括号张成的李代数维数等于位形空间维数 $n$（即\emph{LARC 成立}），则系统\textbf{小时间局部可控}（small-time locally controllable, STLC）——从任意位形出发，可在任意短时间内到达其邻域内任意位形。
\end{theorem}

对差速底盘，把 \cref{eq:mw-diff-bracket} 生成的 $\mathbf{g}_3=[\mathbf{g}_1,\mathbf{g}_2]=(\sin\theta,-\cos\theta,0)^\top$ 与原两场并列，检查秩：
\begin{equation}
  \operatorname{rank}\begin{pmatrix}\cos\theta&0&\sin\theta\\\sin\theta&0&-\cos\theta\\0&1&0\end{pmatrix}=3,\qquad
  \det=-1\cdot\big(\cos\theta\cdot(-\cos\theta)-\sin\theta\cdot\sin\theta\big)=\cos^2\theta+\sin^2\theta=1\neq0,
  \label{eq:mw-larc}
\end{equation}
处处满秩。即：$\mathbf{g}_1,\mathbf{g}_2$ 只张成 $2$ 维分布（不含横向），但加上一阶李括号 $\mathbf{g}_3$ 后张满整个 $\SE(2)$ 的 $3$ 维。LARC 成立 $\Rightarrow$ 差速车可从\emph{任意}位形到达\emph{任意}位形。

\begin{table}[H]\centering\small
\caption{Frobenius 与 Chow：两个判据回答两个不同问题}
\begin{tabular}{llll}
\toprule
定理 & 回答的问题 & 条件 & 结论 \\
\midrule
Frobenius & 约束是否可积？ & 李括号落在分布内 & 约束完整（可积） \\
Chow--Rashevsky & 系统是否可控？ & 李括号张满全空间 & 系统完全可控 \\
\bottomrule
\end{tabular}
\end{table}

\begin{insight}[同一个李括号，两种解读]\label{ins:mw-twoside}
有趣的是 Frobenius 与 Chow 用的是\emph{同一个}李括号 $[\mathbf{g}_1,\mathbf{g}_2]$，结论却互补。Frobenius 问它\emph{在不在}分布内：在，则约束可积、运动被困在子流形；Chow 问它\emph{补没补齐}全空间：补齐，则系统可控。对差速车，$[\mathbf{g}_1,\mathbf{g}_2]$ 逸出 $2$ 维分布（故非完整），又恰好补成 $3$ 维（故可控）——“非完整”与“可控”在同一个向量上达成了统一。
\end{insight}

\subsection{绕路代价：Ball--Box 定理}
\label{sec:mw-ballbox}

LARC 只保证\emph{存在}到达路径，不保证路径短。横向运动靠李括号“凑”出来，效率远低于直接可控的方向。Sub-Riemannian 几何里的\textbf{Ball--Box 定理}给出量化关系：若生成方向 $\mathbf{d}$ 需要 $k$ 阶李括号，则沿 $\mathbf{d}$ 移动距离 $\varepsilon$ 的最短路径长约为
\begin{equation}
  \ell(\varepsilon)\sim O\!\left(\varepsilon^{1/k}\right).
  \label{eq:mw-ballbox}
\end{equation}

\begin{table}[H]\centering\small
\caption{李括号阶数决定绕路代价}
\begin{tabular}{lcll}
\toprule
系统 & LARC 所需阶数 & 横移 $\varepsilon$ 的路径长 & 规划难度 \\
\midrule
全向底盘 & $0$（无需李括号） & $\varepsilon$（直线） & 低 \\
差速底盘 & $1$ & $O(\sqrt{\varepsilon})$ & 中 \\
阿克曼底盘 & $1$（含最小转弯半径） & $O(\sqrt{\varepsilon})$，下界更大 & 高 \\
拖车系统 & $2$ & $O(\varepsilon^{1/3})$ & 很高 \\
\bottomrule
\end{tabular}
\end{table}

\noindent 这把一条直觉事实量化了：差速车倒库比全向底盘难得多（要反复进退），拖着拖车倒库更是难上加难（要更多更细的机动）。阶数越高，“绕路”越长，规划越棘手。它对规划器选型有直接指导：完整系统（全向）可用 A$^\star$/RRT 任意路径搜索；非完整且 LARC 成立的系统须用 Hybrid A$^\star$、RRT + steering 或 state-lattice 这类尊重运动学的规划器。

\begin{pitfall}[概念误区：Chow 定理保证“快速到达”]
新手以为“差速车满足 LARC，所以能快速到任意位置”。Chow 定理只保证\emph{存在}到达路径，不保证路径长度或时间。靠李括号凑出的横移要反复机动，效率远低于直接方向：差速车横移 $1$\,m 的路径长远大于前进 $1$\,m。路径效率与李括号阶数挂钩——阶数越高，绕路越长。
\end{pitfall}

\begin{practice}
\begin{enumerate}
  \item 对阿克曼底盘算出所需的全部李括号，验证 LARC 是否成立；若成立，需要几阶？
  \item 差速车拖一节被动拖车，$\mathbf{q}=(x,y,\theta_1,\theta_2)$（$\theta_1$ 为拖车朝向）。写出控制向量场，验证 LARC 是否成立、需几阶李括号。
  \item （开放题）若某系统 LARC 需 $k$ 阶李括号，直觉上“绕路”到达侧向位置的路径长约是直线距离的多少倍？（提示：Ball--Box 定理 \cref{eq:mw-ballbox}。）
\end{enumerate}
\end{practice}

% ════════════════════════════════════════════════════════════════
\section{控制向量场视角：漂移自由仿射系统}
\label{sec:mw-affine}

把前面四种底盘的运动方程统一抽象，会发现它们都长成同一副模样：位形的时间导数是控制输入的\emph{线性组合}，且不含与输入无关的“漂移项”。这就是\textbf{漂移自由仿射控制系统}的标准形：
\begin{equation}
  \dot{\mathbf{q}}=\sum_{i=1}^{m}u_i\,\mathbf{g}_i(\mathbf{q})=\mathbf{G}(\mathbf{q})\,\mathbf{u},
  \qquad \mathbf{G}(\mathbf{q})=[\mathbf{g}_1(\mathbf{q})\ \cdots\ \mathbf{g}_m(\mathbf{q})],
  \label{eq:mw-affine}
\end{equation}
其中列向量 $\mathbf{g}_i$ 张成允许速度分布 $\mathcal{D}(\mathbf{q})=\ker\mathbf{A}(\mathbf{q})$，$\mathbf{u}=(u_1,\dots,u_m)^\top$ 是降维后的控制输入，$m=n-\operatorname{rank}\mathbf{A}$。差速底盘的 $\mathbf{G}=[\mathbf{g}_1\ \mathbf{g}_2]$（\cref{eq:mw-diff-vf}）、$\mathbf{u}=(v,\omega)^\top$；阿克曼的 $\mathbf{G}$ 与 $\mathbf{u}=(v,\dot\delta)^\top$ 见 \cref{eq:mw-ack-vf}；全向底盘 $\mathbf{A}=\varnothing$ 时 $\mathbf{G}=\mathbf{I}_3$、$\mathbf{u}=(v_x,v_y,\omega)$。

\begin{insight}[“漂移自由”的两重意义]\label{ins:mw-driftfree}
\cref{eq:mw-affine} 没有形如 $\mathbf{f}_0(\mathbf{q})$ 的漂移项，意味着：输入全为零时机器人静止（无惯性滑行），且运动方向完全由当下输入张成的 $\mathcal{D}(\mathbf{q})$ 决定。这带来两个好处：其一，正运动学就是把 $\mathbf{u}$ 经 $\mathbf{G}(\mathbf{q})$ 投到 $\dot{\mathbf{q}}$，规划器只需统一接口 $\dot{\mathbf{q}}=\mathbf{G}(\mathbf{q})\mathbf{u}$，换底盘只换 $\mathbf{G}$；其二，Chow 定理（\cref{thm:mw-chow}）正是为这类系统量身定制的可控性判据，本章一切“绕路可达”的结论都挂靠在此。注意这是\emph{运动学}模型；一旦把质量、滑移惯性纳入，就会出现漂移项，进入\emph{动力学}范畴（本部后续与 \cref{ch:constrained-dynamics}）。
\end{insight}

\noindent 把约束零空间显式参数化为 $\mathbf{G}(\mathbf{q})$，正是“速度命令投影到可行空间”的工程操作之源：给定一个可能违反约束的速度命令 $\mathbf{v}_{\text{cmd}}$，其在允许速度分布上的正交投影为
\begin{equation}
  \mathbf{v}_{\text{proj}}=(\mathbf{I}-\mathbf{A}^{\dagger}\mathbf{A})\,\mathbf{v}_{\text{cmd}},
  \label{eq:mw-project}
\end{equation}
$\mathbf{A}^\dagger$ 为 $\mathbf{A}$ 的 Moore--Penrose 伪逆，$\mathbf{I}-\mathbf{A}^\dagger\mathbf{A}$ 恰是到 $\ker\mathbf{A}(\mathbf{q})=\mathcal{D}(\mathbf{q})$ 的正交投影。差速底盘代入即把任意 $(v_x,v_y,\omega)$ 的横向分量抹掉、保留前向与转动——这与零空间投影在冗余机械臂里的用法（\cref{sec:umm-nullspace}）是同一套数学。

% ════════════════════════════════════════════════════════════════
\section{轮胎滑移与约束松弛：当纯滚动假设失效}
\label{sec:mw-slip}

\paragraph{纯滚动何时失效。}
前几节的推导都压在一个假设上：\emph{纯滚动 $+$ 不横滑}。它在低速（$<1$\,m/s）、干硬地面、轻载、直线运动时近似成立；但在湿滑地砖、松软沙土、急加速/急刹、高速过弯时严重失效。本节量化“滑了多少”，并把理想约束 $\mathbf{A}\dot{\mathbf{q}}=\mathbf{0}$ 松弛为带残差的 $\mathbf{A}\dot{\mathbf{q}}=\mathbf{r}(\kappa,\alpha)$。

\begin{insight}[滑移 = Pfaffian 约束的残差]\label{ins:mw-residual}
$\mathbf{A}(\mathbf{q})\dot{\mathbf{q}}=\mathbf{0}$ 描述的是\emph{理想运动学}——假设轮地刚性接触、零滑动。真实轮胎是弹性体，接触区有有限面积，轮胎力靠变形产生而非理想约束。滑移建模就是在 Pfaffian 框架里引入“约束的偏差”：从 $\mathbf{A}\dot{\mathbf{q}}=\mathbf{0}$ 变为
\begin{equation}
  \mathbf{A}(\mathbf{q})\dot{\mathbf{q}}=\mathbf{r}(\kappa,\alpha),
  \label{eq:mw-A-residual}
\end{equation}
残差 $\mathbf{r}$ 由滑移参数 $\kappa,\alpha$ 决定，$\kappa=\alpha=0$ 时退回理想纯滚动。
\end{insight}

\paragraph{两个滑移参数。}
滑移用两个无量纲量刻画。\emph{纵向滑移比} $\kappa$：
\begin{equation}
  \kappa=\frac{r\omega-v_x}{\max(|v_x|,\,v_{\text{low}})},
  \label{eq:mw-kappa}
\end{equation}
$\kappa=0$ 为纯滚动（$r\omega=v_x$），$\kappa>0$ 为驱动打滑（轮缘快于地面，如起步空转），$\kappa<0$ 为制动打滑（地面快于轮缘，如急刹），$\kappa=1$ 为完全空转（$v_x=0,\omega\neq0$），$\kappa=-1$ 为完全抱死（$\omega=0,v_x\neq0$）。分母里的 $v_{\text{low}}$（典型 $0.05$–$0.1$\,m/s）是\emph{低速保护}：低速时分母趋零会让 $\kappa$ 剧烈震荡、失去物理意义。\emph{侧向滑移角} $\alpha$：
\begin{equation}
  \alpha=\arctan\!\left(\frac{v_y}{|v_x|+v_{\text{low}}}\right),
  \label{eq:mw-alpha}
\end{equation}
$\alpha=0$ 无侧滑，$\alpha>0$ 向左侧滑、$\alpha<0$ 向右侧滑。

\paragraph{Pacejka 魔术公式。}
Hans Pacejka 提出的“魔术公式”（Magic Formula）是汽车工业最广用的轮胎力模型。称“魔术”是因为它\emph{没有直接的物理推导}，纯属经验拟合，却惊人地适配各种轮胎与工况。其纯纵向或纯侧向的简化形式为
\begin{equation}
  F=F_z\,D\,\sin\!\Big(C\,\arctan\!\big(B\,s-E\,(B\,s-\arctan(B\,s))\big)\Big),
  \label{eq:mw-pacejka}
\end{equation}
$s$ 是滑移参数（纵向取 $\kappa$、侧向取 $\alpha$），$B,C,D,E$ 是四个经验参数（见 \cref{tab:mw-pacejka}）。无论轮胎模型多复杂，接触力终究受\emph{摩擦锥}约束
\begin{equation}
  \sqrt{F_x^2+F_y^2}\le\mu F_z,
  \label{eq:mw-friction-cone}
\end{equation}
这是一切轮式接触约束的最终边界，与足式接触力的摩擦锥（\cref{eq:cd-signorini-foot} 所在的接触约束族）同源。

\begin{table}[H]\centering\small
\caption{Pacejka 魔术公式的四个参数}
\label{tab:mw-pacejka}
\begin{tabular}{lll}
\toprule
参数 & 名称 & 物理含义 \\
\midrule
$B$ & 刚度因子 & 小滑移时力--滑移斜率 \\
$C$ & 形状因子 & 曲线形状（纵向约 $1.65$，侧向约 $1.3$） \\
$D$ & 峰值因子 & 最大摩擦系数 $\mu_{\max}$ \\
$E$ & 曲率因子 & 峰值后的下降速率 \\
\bottomrule
\end{tabular}
\end{table}

\begin{figure}[ht]
\centering
\begin{tikzpicture}[>=Stealth, scale=1.0]
  \draw[->] (0,0) -- (5.4,0) node[right, font=\small]{滑移 $s$};
  \draw[->] (0,0) -- (0,3.0) node[above, font=\small]{$F/F_z$};
  % magic-formula-like curve: rise, peak, slight fall (control points, no trig units)
  \draw[very thick, main!70!black]
    plot[smooth, tension=0.7]
    coordinates {(0,0) (0.5,1.35) (1.0,2.25) (1.55,2.55) (2.3,2.4) (3.4,2.18) (5.0,2.05)};
  % peak marker
  \draw[dashed, gray!60!black] (0,2.55) -- (1.55,2.55) -- (1.55,0);
  \node[gray!60!black, font=\scriptsize, left] at (0,2.55) {$\mu_{\max}$};
  % region labels
  \draw[<->, second!70!black] (0,-0.55) -- (0.95,-0.55);
  \node[second!70!black, font=\scriptsize] at (0.5,-0.85) {线性区};
  \draw[<->, third!70!black] (0.95,-0.55) -- (2.3,-0.55);
  \node[third!70!black, font=\scriptsize] at (1.6,-0.85) {峰值区};
  \draw[<->, red!60!black] (2.3,-0.55) -- (5.0,-0.55);
  \node[red!60!black, font=\scriptsize] at (3.6,-0.85) {滑动区};
  % linear tangent at origin
  \draw[dotted, second!70!black] (0,0) -- (1.2,2.85);
\end{tikzpicture}
\caption{Pacejka 魔术公式 \cref{eq:mw-pacejka} 的典型 S 形力--滑移曲线：小滑移时力近似线性增长（斜率 $\propto B$），到峰值区达最大摩擦 $\mu_{\max}F_z$，进入滑动区后力反而下降。线性区可用 $F\approx C_s\,s$ 近似。}
\label{fig:mw-pacejka}
\end{figure}

\begin{table}[H]\centering\small
\caption{Pacejka 曲线的三个特征区}
\begin{tabular}{llll}
\toprule
区域 & 滑移范围 & 物理行为 & 控制意义 \\
\midrule
线性区 & $|s|<0.03$ & 力--滑移近似线性 $F\approx C_s s$ & 正常操作区，PD 控制有效 \\
峰值区 & $|s|\approx0.05$–$0.1$ & 摩擦达最大 $F_{\max}=DF_z$ & 最大制/驱动力，ABS/TCS 目标 \\
滑动区 & $|s|>0.15$ & 摩擦力下降（$F<F_{\max}$） & 失控区，须减速或切换模式 \\
\bottomrule
\end{tabular}
\end{table}

\paragraph{轮足/小轮场景：线性模型足够。}
对轮足机器人或小轮（半径约 $0.05$\,m），接触面积小、轮胎力学与汽车差异大，通常\emph{不必}精拟 Pacejka 参数——但力--滑移曲线的\emph{定性形状}仍适用（小滑移线性、大滑移饱和）。控制器设计常用线性轮胎模型即可：
\begin{equation}
  F_x=C_\kappa\,\kappa,\qquad F_y=C_\alpha\,\alpha,
  \label{eq:mw-linear-tire}
\end{equation}
$C_\kappa$ 为纵向刚度（滑移比 $\kappa$ 无量纲，故 $C_\kappa$ 单位为 N；小轮典型 $500$–$2000$\,N \pz{量级依轮径/材质/法向载荷而异，乘用车整轮的纵向刚度可达数万 N}）、$C_\alpha$ 为侧偏刚度（$\alpha$ 为弧度，故单位 N/rad；小轮典型 $300$–$1500$\,N/rad \pz{同样依平台而异：乘用车整轮 $C_\alpha$ 量级在 $10^4$–$10^5$\,N/rad，此处取值面向轻载小轮}），在 $|\kappa|<0.05$、$|\alpha|<3^\circ$ 内精度足够。

\paragraph{残差与里程计漂移。}
把滑移接回约束残差 \cref{eq:mw-A-residual}：当存在滑移时 $\mathbf{A}(\mathbf{q})\dot{\mathbf{q}}=\mathbf{r}(\kappa,\alpha)\neq\mathbf{0}$。线性/小角近似下，纵向残差 $r_{\text{lon}}\approx r\dot\theta\cdot\kappa$、侧向残差 $r_{\text{lat}}\approx v\sin\alpha$。$\kappa=\alpha=0$ 时残差归零，回到理想纯滚动；滑移越大，残差越大，纯滚动运动学的预测越不准。基于纯滚动假设的轮式里程计（wheel odometry）因此会积累定位误差，误差增长率近似
\begin{equation}
  \dot{e}_{\text{pos}}\approx v\,\sqrt{\kappa^2+\alpha^2}.
  \label{eq:mw-odom-drift}
\end{equation}
取 $v=1$\,m/s、$\kappa=0.05$、$\alpha=3^\circ$，得 $\dot{e}_{\text{pos}}\approx0.07$\,m/s——每秒漂移约 $7$\,cm，$10$\,秒就累积约 $70$\,cm，对“在门口精确放包裹”的配送机器人完全不可接受。这正是实际导航不能只靠轮式里程计、必须融合 IMU/视觉/LiDAR 的原因：滑移大时调高轮速观测的协方差、调低对里程计的信任，是约束残差观点给状态估计的直接指导。

\begin{pitfall}[编程陷阱：低速时直接套滑移比公式]
直接写 $\kappa=(r\omega-v_x)/v_x$、不加低速保护：机器人低速启动时分母趋零，$\kappa$ 突然飙到 $\pm100$，触发安全保护或令控制器震荡。正确做法：分母用 $\max(|v_x|,v_{\text{low}})$，并在 $|v_x|<v_{\text{low}}$ 时改用速度残差 $r\omega-v_x$（绝对量）而非比值 $\kappa$。
\end{pitfall}

\begin{pitfall}[概念误区：轮足机器人必须用 Pacejka 模型]
Pacejka 模型有 $6$–$20$ 个经验参数，要大量轮胎测试数据标定。轮足机器人的“轮胎”多是橡胶包覆的小轮，与汽车轮胎差异巨大。工程上轮足更常用\emph{线性力--滑移模型}（\cref{eq:mw-linear-tire}）或\emph{简单摩擦锥}（\cref{eq:mw-friction-cone}），因为参数少、可标定、算得快。按精度需求选模型复杂度：户外高速可考虑简化 Pacejka，室内低速线性模型足矣。
\end{pitfall}

\begin{practice}
\begin{enumerate}
  \item 实现带低速保护的纵向滑移比函数（输入 $v_{\text{wheel}},v_{\text{ground}},v_{\text{low}}$），对 $v_{\text{ground}}\in\{0.01,0.1,1.0\}$、$v_{\text{wheel}}\in\{0,0.5,1.0,2.0\}$ 各组合画热力图。
  \item 用 Pacejka 简化公式（取 $B=10,C=1.5,D=0.8,E=-0.5$）画侧向力 vs 滑移角曲线，标出峰值摩擦对应的滑移角，并解释峰值为何不在最大滑移角处。
  \item 设计“滑移感知权重调度器”：输入四轮滑移比，输出四个权重 $w_i\in[0,1]$ 调节状态估计对各轮速观测的信任，要求 $|\kappa_i|<0.05$ 时 $w_i=1$、$|\kappa_i|>0.3$ 时 $w_i=0$、中间平滑过渡。
\end{enumerate}
\end{practice}

% ════════════════════════════════════════════════════════════════
\section{承上启下}
\label{sec:mw-bridge}

本章把“轮子与地面的接触”翻译成了速度层的 Pfaffian 约束，并追问了它的两面：可积性（Frobenius，复用 \cref{thm:cd-frobenius}）与可控性（Chow，本章补全），最后用约束残差把理想纯滚动与真实滑移连了起来。这套\emph{速度约束 $\to$ 控制向量场 $\to$ 可控性}的链条，是移动操作一切运动学的地基。

下一章 \cref{ch:mm-kinematics} 将沿两个方向推进：其一，把本章 $\SE(2)$ 底盘的速度 $\mathbf{v}^b$ 与机械臂雅可比拼成移动操作的\emph{联合雅可比} $\mathbf{V}_{ee}=[\mathbf{J}_b\ \mathbf{J}_a]\boldsymbol{\nu}$，其中底盘块 $\mathbf{J}_b$ 的列正是本章 $\mathbf{G}(\mathbf{q})$（允许速度分布的基）；其二，处理底盘非完整约束与机械臂冗余（\cref{sec:ak-redundancy}）叠加后的零空间分配与奇异性（\cref{sec:ak-singularity}）。而本章把底盘写成漂移自由仿射系统 $\dot{\mathbf{q}}=\mathbf{G}(\mathbf{q})\mathbf{u}$ 的标准形，也正是后续移动操作最优控制（轨迹优化把它当作动力学约束）的接口。

% ════════════════════════════════════════════════════════════════
\section*{本章小结}

\begin{table}[H]\centering\small
\caption{符号表}
\begin{tabular}{lll}
\toprule
符号 & 含义 & 首次出现 \\
\midrule
$\mathbf{q}=(x,y,\theta)$ & 底盘位形（$\SE(2)$） & \cref{eq:mw-se2} \\
$\mathbf{v}^b=(v_x^b,v_y^b,\omega)$ & 体坐标系速度（前向/左向/偏航） & \cref{eq:mw-body2world} \\
$\mathbf{R}(\theta)$ & 平面旋转矩阵（体$\to$世界） & \cref{eq:mw-body2world} \\
$\mathbf{A}(\mathbf{q})$ & Pfaffian 约束矩阵 & \cref{eq:mw-pfaffian} \\
$\mathcal{D}(\mathbf{q})=\ker\mathbf{A}$ & 允许速度分布 & \cref{eq:mw-distribution} \\
$\mathbf{g}_i(\mathbf{q}),\ u_i$ & 控制向量场 / 输入 & \cref{eq:mw-disk-vf} \\
$[\mathbf{g}_i,\mathbf{g}_j]$ & 李括号 & \cref{eq:mw-involutive} \\
$\mathbf{G}(\mathbf{q})=[\mathbf{g}_1\,\cdots\,\mathbf{g}_m]$ & 仿射系统输入矩阵 & \cref{eq:mw-affine} \\
$r,\ L,\ L_{\text{wb}},\ \delta$ & 轮半径/轮距/轴距/转向角 & \cref{eq:mw-diff-fk,eq:mw-ack-q} \\
$\kappa,\ \alpha$ & 纵向滑移比 / 侧向滑移角 & \cref{eq:mw-kappa,eq:mw-alpha} \\
$B,C,D,E$ & Pacejka 四参数 & \cref{eq:mw-pacejka} \\
$\mathbf{r}(\kappa,\alpha)$ & 约束残差（滑移） & \cref{eq:mw-A-residual} \\
$\mu,\ F_z$ & 摩擦系数 / 法向力 & \cref{eq:mw-friction-cone} \\
\bottomrule
\end{tabular}
\end{table}

\begin{table}[H]\centering\small
\caption{定理 / 公式速查表}
\begin{tabular}{lll}
\toprule
定理 / 公式 & 一句话说明 & 对应 \\
\midrule
体$\to$世界速度 \cref{eq:mw-body2world} & 旋转矩阵搬运速度，约束推导的坐标基建 & \cref{sec:mw-se2} \\
Pfaffian 形式 \cref{eq:mw-pfaffian} & $\mathbf{A}\dot{\mathbf{q}}=\mathbf{0}$，零空间即瞬时自由度 & \cref{eq:mw-distribution} \\
四底盘 $\mathbf{A}(\mathbf{q})$ & 三步框架推导，秩读出自由度 & \cref{tab:mw-chassis} \\
Frobenius（专化） & $[\mathbf{g}_i,\mathbf{g}_j]\in\mathcal{D}$ $\Leftrightarrow$ 可积 & \cref{thm:mw-frobenius} \\
差速李括号 \cref{eq:mw-diff-bracket} & $[\mathbf{g}_1,\mathbf{g}_2]$ = 侧向 = 泊车本质 & \cref{ins:mw-parking} \\
Chow / LARC \cref{eq:mw-larc} & 李括号张满 $\Rightarrow$ 可控（绕路可达） & \cref{thm:mw-chow} \\
Ball--Box \cref{eq:mw-ballbox} & $k$ 阶括号 $\Rightarrow$ 绕路 $O(\varepsilon^{1/k})$ & \cref{sec:mw-ballbox} \\
仿射系统 \cref{eq:mw-affine} & $\dot{\mathbf{q}}=\mathbf{G}(\mathbf{q})\mathbf{u}$，漂移自由标准形 & \cref{sec:mw-affine} \\
滑移残差 \cref{eq:mw-A-residual} & $\mathbf{A}\dot{\mathbf{q}}=\mathbf{r}(\kappa,\alpha)$，纯滚动松弛 & \cref{ins:mw-residual} \\
Pacejka \cref{eq:mw-pacejka} & 经验力--滑移 S 曲线，四参数 & \cref{fig:mw-pacejka} \\
里程计漂移 \cref{eq:mw-odom-drift} & $\dot e\approx v\sqrt{\kappa^2+\alpha^2}$，融合动机 & \cref{sec:mw-slip} \\
\bottomrule
\end{tabular}
\end{table}

\begin{table}[H]\centering\small
\caption{知识点总表}
\begin{tabular}{clll}
\toprule
\# & 知识点 & 核心要点 & 难度 \\
\midrule
1 & 完整 vs 非完整 & 前者降位形维、后者锁瞬时速度方向 & $\bigstar\bigstar$ \\
2 & $\SE(2)$ 运动学 & 体$\leftrightarrow$世界速度，$n_q=n_v=3$ & $\bigstar\bigstar$ \\
3 & Pfaffian 推导 & 三步框架，四底盘 $\mathbf{A}(\mathbf{q})$ & $\bigstar\bigstar\bigstar$ \\
4 & Frobenius 可积性 & 李括号在不在分布内 & $\bigstar\bigstar\bigstar$ \\
5 & Chow / LARC & 李括号张满 $\Rightarrow$ 可控 & $\bigstar\bigstar\bigstar$ \\
6 & Ball--Box 绕路代价 & $O(\varepsilon^{1/k})$，阶数定难度 & $\bigstar\bigstar\bigstar\bigstar$ \\
7 & 仿射控制系统 & $\dot{\mathbf{q}}=\mathbf{G}(\mathbf{q})\mathbf{u}$ 标准形 & $\bigstar\bigstar\bigstar$ \\
8 & 滑移建模 & $\kappa,\alpha$、Pacejka、摩擦锥 & $\bigstar\bigstar\bigstar$ \\
9 & 约束残差与漂移 & $\mathbf{A}\dot{\mathbf{q}}=\mathbf{r}$，里程计融合 & $\bigstar\bigstar\bigstar$ \\
\bottomrule
\end{tabular}
\end{table}

\section*{延伸阅读}
\begin{itemize}
  \item \textbf{非完整运动规划的系统化源头}：Murray、Li 与 Sastry《A Mathematical Introduction to Robotic Manipulation》\cite{murray1994mathematical}——李括号、Chow 定理、非完整 steering 的经典论述（本章 \cref{sec:mw-integrability,sec:mw-controllability} 的主要依据）。
  \item \textbf{采样规划与非完整系统}：LaValle《Planning Algorithms》\cite{lavalle2006planning}——把 RRT 等推广到非完整系统、运动原语与 lattice 规划。
  \item \textbf{约束力学与可积性判据}：本书 \cref{ch:constrained-dynamics}（Pfaffian 统一形式、Frobenius/Signorini 判型、DAE 与接触）与 \cref{ch:geometric-mechanics}（完整/非完整的几何定义 \cref{def:gm-constraint}）——本章的一般理论出处。
  \item \textbf{现代运动学教材视角}：Lynch 与 Park《Modern Robotics》\cite{lynch2017modern}——$\SE(2)$/$\SE(3)$ 旋量理论与轮式机器人运动学的统一处理。
  \item \textbf{轮胎力学}：Pacejka《Tire and Vehicle Dynamics》\cite{pacejka2012tire}——魔术公式 \cref{eq:mw-pacejka} 的权威出处与参数标定；非完整力学几何见 Bloch《Nonholonomic Mechanics and Control》\cite{bloch2015nonholonomic}。
\end{itemize}

\section*{本章与后续章节的关系}
\begin{table}[H]\centering\small
\begin{tabular}{lll}
\toprule
后续章节 & 关系 & 本章铺垫 \\
\midrule
\cref{ch:mm-kinematics} 移动操作运动学 & 底盘 $+$ 臂联合雅可比 & $\mathbf{G}(\mathbf{q})$ 即底盘雅可比块 $\mathbf{J}_b$ \\
\cref{ch:constrained-dynamics} 约束动力学 & 把约束写进动力学 KKT & Pfaffian 形式、Frobenius 判据 \\
\cref{ch:unified-multimodal-mpc} 统一多模态 MPC & 接触/滚动约束的统一处理 & 接触=滚动的残差松弛视角 \\
\cref{sec:umm-nullspace} 零空间投影 & 命令投影到可行子空间 & \cref{eq:mw-project} 的零空间投影 \\
\bottomrule
\end{tabular}
\end{table}

\section*{故障排查手册}
\begin{enumerate}
  \item \textbf{症状}：差速/阿克曼机器人跟踪规划路径时\emph{频繁偏离、原地打转、轮胎打滑异响}。\textbf{排查}：规划器是否用了 A$^\star$/Dijkstra 等\emph{完整系统}搜索（\cref{sec:mw-motivation} 反面）？应换成 Hybrid A$^\star$/state-lattice 等尊重非完整约束的规划器。
  \item \textbf{症状}：约束矩阵\emph{只在 $\theta=0$ 时正确}、其他朝向运动方向错乱。\textbf{原因}：在世界系写了常数 $\mathbf{A}$、忘了它依赖 $\theta$（\cref{sec:mw-se2} 陷阱）。\textbf{对策}：先在体系写约束，再经 \cref{eq:mw-body2world-comp} 变到世界系。
  \item \textbf{症状}：数值李括号与解析结果\emph{差异大}、$\theta$ 近 $0$ 或 $\pi$ 时尤甚。\textbf{原因}：有限差分步长过大（\cref{sec:mw-diff-bracket} 陷阱）。\textbf{对策}：改用符号微分或自动微分。
  \item \textbf{症状}：机器人低速启动时滑移比\emph{突然飙到 $\pm100$}、触发保护或控制震荡。\textbf{原因}：滑移比分母趋零（\cref{sec:mw-slip} 陷阱）。\textbf{对策}：分母用 $\max(|v_x|,v_{\text{low}})$，极低速时改用速度残差而非比值。
  \item \textbf{症状}：纯轮式里程计\emph{定位持续漂移}，湿滑/松软地面尤甚。\textbf{原因}：滑移使 $\mathbf{A}\dot{\mathbf{q}}=\mathbf{r}\neq\mathbf{0}$、纯滚动假设失效（\cref{eq:mw-odom-drift}）。\textbf{对策}：融合 IMU/视觉/LiDAR，按估计滑移调高轮速观测协方差、调低里程计信任。
  \item \textbf{症状}：多轮底盘的约束系统\emph{秩亏、自相矛盾}。\textbf{原因}：把各轮约束简单叠加、忽略轮间几何耦合（\cref{tab:mw-chassis} 后陷阱）。\textbf{对策}：堆叠全部约束后分析 $\mathbf{A}(\mathbf{q})$ 的秩，只保留独立行。
\end{enumerate}
